\section{Block Log-Det Condensation}
\label{sec:condensation}

This section studies the log-det center of an SDP whose variable is a tuple
of PSD blocks with one shared trace constraint $\sum_g\tr Z_g=1$.  Suppose
one block, the winner, has a cost whose least eigenvalue is smaller than
that of every other block.  As the path parameter decreases, the trace
multiplier approaches this eigenvalue, where the winner's resolvent has a
pole.  The losing blocks can hold only a bounded amount of trace near that
eigenvalue, so the remaining trace condenses on the winner.  We show that
the winner's trace mass controls three quantities: the condition number of
the equality-reduced Hessian, the trace distance
between the centers with and without a winner (the visibility of the
winner), and the query cost of preparing the center as a quantum state or
extracting the winner from it.

We work with an abstract two-cost block model.  \Cref{sec:winner-take-all}
gives a sparse SDP that realizes it, and \cref{sec:state-conversion} proves
a stronger lower bound for mixed-state output.  In this section only, $G$ is
the number of blocks, $d$ their common dimension, $m$ the multiplicity of
the least eigenvalue of the winning cost (not the number of equality
constraints), $\tau$ the path parameter (called $\mu$ in
\cref{def:central-path}), $\alpha=\tau G$, and $p$ the trace of the winning
block (the winner mass).  The $C$- and $\Gamma$-symbols and $d,b,c$ are also
local, and the resolvent trace $A(x)$ is unrelated to the constraint matrix
$A$ of \cref{sec:prelim-conic}.
\newcommand{\condtag}[1]{\tag{#1}}

\subsection{Two-cost model and exact center}
\label{sec:condensation-model}

Fix real symmetric costs $C_*,C_L\in\Sym^d$.  There is one winning block
with cost $C_*$ and $G-1$ identical losing blocks with cost $C_L$.  After
subtracting the winning spectral minimum $v_*$ from both costs, suppose
\begin{align}
 \operatorname{spec}(C_*-v_*I)
  &=\{\underbrace{0,\ldots,0}_{m},\gamma_{m+1},\ldots,\gamma_d\},
   &\gamma_i&>0,                                             \label{eq:cond-winner-spectrum}\\
 \operatorname{spec}(C_L-v_*I)
  &=\{\delta_1,\ldots,\delta_d\},
   &\delta_j&>0.                                             \label{eq:cond-loser-spectrum}
\end{align}
Here $1\le m\le d$, and the winning block is unique.  Unless otherwise
stated, $d$, $m$, and all positive gaps are fixed while $G$ grows.

At path parameter $\tau>0$, consider
\begin{equation}
 \min_{\substack{Z_g\succ0\\ \sum_g\tr Z_g=1}}
 \left\{\inner{C_*}{Z_*}+\sum_{g=2}^G\inner{C_L}{Z_g}
             -\tau\sum_{g=1}^G\log\det Z_g\right\}.
 \label{eq:cond-center-program}
\end{equation}
For $x>0$, define the winner and loser resolvent traces
\begin{equation}
 A(x)=\frac{m}{x}+\sum_{i=m+1}^d\frac1{x+\gamma_i},
 \qquad
 B(x)=\sum_{j=1}^d\frac1{x+\delta_j}.
 \label{eq:cond-resolvents}
\end{equation}
An empty regular sum is zero.  Also put
\begin{equation}
 B_0=B(0)=\sum_{j=1}^d\delta_j^{-1},
 \qquad
 \beta=-B'(0)=\sum_{j=1}^d\delta_j^{-2},
 \qquad
 \alpha_*=B_0^{-1}.
 \label{eq:cond-capacity}
\end{equation}

\begin{proposition}[Exact block center and finite concentration bounds]
\label{prop:cond-exact-center}
The unique minimizer of \eqref{eq:cond-center-program} is
\begin{equation}
 Z_*=\tau(C_*-\lambda I)^{-1},\qquad
 Z_g=\tau(C_L-\lambda I)^{-1}\quad(g>1),
 \label{eq:cond-exact-center}
\end{equation}
where $\lambda=v_*-x$ and $x$ is the unique positive solution of
\begin{equation}
 1=\tau\{A(x)+(G-1)B(x)\}.
 \label{eq:cond-secular}
\end{equation}
The winner mass $p_G=\tr Z_*=\tau A(x)$ satisfies, with
$\delta_{\max}=\max_j\delta_j$,
\begin{align}
 p_G&\ge1-\tau(G-1)B_0,                                  \label{eq:cond-finite-lower}\\
 p_G\ge p_0>0&\quad\Longrightarrow\quad
 x\le\frac{d\tau}{p_0},\qquad
 1-p_G\ge\frac{d\tau(G-1)}{\delta_{\max}+d\tau/p_0}.
                                                               \label{eq:cond-finite-converse}
\end{align}
Consequently, for $G\ge2$ and $0<\varepsilon<1$,
\begin{align}
 \tau\le\frac{\varepsilon}{B_0(G-1)}
 &\quad\Longrightarrow\quad p_G\ge1-\varepsilon,          \label{eq:cond-epsilon-sufficient}\\
 p_G\ge1-\varepsilon
 &\quad\Longrightarrow\quad
 \tau\le\frac{\varepsilon\delta_{\max}}
 {d\{G-1-\varepsilon/(1-\varepsilon)\}},                  \label{eq:cond-epsilon-necessary}
\end{align}
whenever the denominator in \eqref{eq:cond-epsilon-necessary} is positive.
For fixed $G$ and $\tau\downarrow0$, if
$A_{\rm reg}(0)=\sum_{i=m+1}^d\gamma_i^{-1}$, then
\begin{align}
 x&=m\tau+m\{A_{\rm reg}(0)+(G-1)B_0\}\tau^2+O_G(\tau^3),
                                                               \label{eq:cond-small-tau-x}\\
 p_G&=1-(G-1)B_0\tau+m(G-1)\beta\tau^2+O_G(\tau^3).
                                                               \label{eq:cond-small-tau-p}
\end{align}
\end{proposition}

\begin{proof}
The objective is strictly convex on the affine slice and diverges at its
relative boundary, so it has a unique minimizer.  With the multiplier sign
from \cref{def:trace-logdet-center}, stationarity gives
$C_g-\lambda I=\tau Z_g^{-1}$.  Positive definiteness requires
$\lambda<v_*$, and taking traces gives \eqref{eq:cond-secular}.  Its
right-hand side is continuous and strictly decreasing from infinity to zero,
which proves existence and uniqueness of $x$.

The loser mass is $1-p_G=\tau(G-1)B(x)$.  Since $B(x)\le B_0$, this proves
\eqref{eq:cond-finite-lower}.  Since $A(x)\le d/x$, $p_G\ge p_0$ implies
$x\le d\tau/p_0$.  The bound
$B(x)\ge d/(x+\delta_{\max})$ then proves
\eqref{eq:cond-finite-converse}.  Equations
\eqref{eq:cond-epsilon-sufficient}--\eqref{eq:cond-epsilon-necessary} follow
by substituting $p_0=1-\varepsilon$ and rearranging.

Finally, Taylor expansion gives
$A(x)=m/x+A_{\rm reg}(0)+O(x)$ and
$B(x)=B_0-\beta x+O(x^2)$.  Substitute the ansatz
$x=m\tau+c\tau^2+O_G(\tau^3)$ into \eqref{eq:cond-secular}.  Subtracting
one, dividing by $\tau$, and taking the constant term gives
$c=m\{A_{\rm reg}(0)+(G-1)B_0\}$.  Substitution into
$1-p_G=\tau(G-1)B(x)$ gives \eqref{eq:cond-small-tau-p}.
\end{proof}

Thus the winner mass stays bounded away from zero only when $\tau=O(1/G)$,
and it is close to one when $\tau$ is a small enough multiple of $1/G$.
The next result resolves this scale.

\subsection{Capacity transition}
\label{sec:condensation-threshold}

With $\tau=\alpha/G$, the losers hold total trace
$\tau(G-1)B(x)<\alpha B_0$.  We call $\alpha_*=1/B_0$ the capacity and the
cases $\alpha<\alpha_*$, $\alpha=\alpha_*$, and $\alpha>\alpha_*$
subcritical, critical, and supercritical.

\begin{theorem}[Block log-det condensation]
\label{thm:condensation-threshold}
Set $\tau=\alpha/G$ for fixed $\alpha>0$.  Under
\eqref{eq:cond-winner-spectrum}--\eqref{eq:cond-loser-spectrum}, the following
limits hold as $G\to\infty$:
\begin{alignat}{3}
 0<\alpha<\alpha_*:&\quad
 &\frac{x}{\tau}&\longrightarrow\frac{m}{1-\alpha B_0},
 &\qquad p_G&\longrightarrow1-\alpha B_0;                 \label{eq:cond-subcritical}\\
 \alpha=\alpha_*:&
 &\sqrt G\,x&\longrightarrow\sqrt{m/\beta},
 &\sqrt G\,p_G&\longrightarrow\alpha_*\sqrt{m\beta};    \label{eq:cond-critical}\\
 \alpha>\alpha_*:&
 &x&\longrightarrow x_\alpha,
 &G p_G&\longrightarrow\alpha A(x_\alpha),               \label{eq:cond-supercritical}
\end{alignat}
where $x_\alpha>0$ is the unique solution of
$B(x_\alpha)=1/\alpha$.
\end{theorem}

\begin{proof}
Write the secular equation as
\begin{equation}
 1=\frac{\alpha}{G}A(x)+\alpha(1-G^{-1})B(x).
 \label{eq:cond-scaled-secular}
\end{equation}
If $\alpha<\alpha_*$, $x$ must tend to zero: along any subsequence bounded
away from zero, the first term vanishes and the second has limit at most
$\alpha B_0<1$.  Using the expansions following
\eqref{eq:cond-small-tau-p} in \eqref{eq:cond-scaled-secular} gives
\[
 1=\frac{\alpha m}{Gx}+\alpha B_0+o(1),
\]
so $Gx\to\alpha m/(1-\alpha B_0)$ and
$p_G=(\alpha/G)A(x)\to1-\alpha B_0$.

At $\alpha=\alpha_*$ the same compactness argument gives $x\to0$.
After cancelling $\alpha_*B_0=1$, the Taylor expansions give
\begin{equation}
 0=\frac{\alpha_*m}{Gx}-\alpha_*\beta x
       +O(G^{-1})+O(x^2),
 \label{eq:cond-critical-balance}
\end{equation}
where the bounded regular part of $A$ is absorbed in $O(G^{-1})$.
Monotonicity first sandwiches $x$ between constant multiples of
$G^{-1/2}$; dividing \eqref{eq:cond-critical-balance} by $G^{-1/2}$ then
yields $\sqrt Gx\to\sqrt{m/\beta}$.  Since
$A(x)=m/x+O(1)$, the mass limit in \eqref{eq:cond-critical} follows.

If $\alpha>\alpha_*$, $x$ cannot tend to zero because the loser term alone
would tend to $\alpha B_0>1$.  It also cannot diverge.  Every convergent
subsequence therefore has a positive limit, at which the winner term in
\eqref{eq:cond-scaled-secular} vanishes and $B(x)=1/\alpha$.  Strict
monotonicity of $B$ makes that limit unique.  Multiplying
$p_G=(\alpha/G)A(x)$ by $G$ proves \eqref{eq:cond-supercritical}.
\end{proof}

Below capacity, the winner's pole term $m/x$ supplies the trace the losers
cannot hold; at capacity, it balances the first-order decrease $\beta x$ of
the loser trace; above capacity, the loser equation alone determines
$x_\alpha$.  The assumption $\delta_j>0$ is essential: a zero
loser gap makes $B_0$ infinite, and there is no finite capacity.

\subsection{Equality-reduced Hessian and conditioning}
\label{sec:condensation-conditioning}

All condition numbers below are those of \cref{def:reduced-kappa}: the
block Hessian restricted to the tangent space of the trace constraint, in
Frobenius coordinates.  They say nothing about preconditioned systems or the
saddle-point KKT matrix.  At the center put
$D_*=C_*-\lambda I$ and $D_L=C_L-\lambda I$.  On the tangent space
\begin{equation}
 \mathcal T=\left\{Y=(Y_1,\ldots,Y_G):
                    \sum_{g=1}^G\tr Y_g=0\right\},
 \label{eq:cond-tangent}
\end{equation}
the barrier-scaled log-det Hessian has quadratic form
\begin{equation}
\mathcal Q(Y)=\frac1\tau\left\{
  \tr(D_*Y_*D_*Y_*)+\sum_{g=2}^G\tr(D_LY_gD_LY_g)\right\}.
 \label{eq:cond-hessian-form}
\end{equation}
This \(\mathcal Q\) is the quadratic form of
\(\tau\nabla^2F_{\log\det}\) at the center, so
\(\kappa(\mathcal Q|_{\mathcal T})=\kappa_{\rm red}\); the factor \(\tau\)
does not change curvature ratios.

Let $\delta_{\min}=\min_j\delta_j$, let $r$ be its multiplicity, and, when
$m<d$, let $\gamma_{\min}=\min_{i>m}\gamma_i$.  Define
\begin{equation}
 b(x)=\min\bigl(\{x+\delta_{\min}\}\cup
                  \{x+\gamma_{\min}:m<d\}\bigr),
 \quad
 \Lambda=\max\bigl(\{\delta_j\}\cup
                  \{\gamma_i:i>m\}\bigr),
 \label{eq:cond-gap-extremes}
\end{equation}
with the second set omitted when $m=d$.

\begin{lemma}[Finite Hessian sandwich]
\label{lem:cond-hessian-sandwich}
For $m\ge2$,
\begin{equation}
 \lambda_{\min}(\mathcal Q|_{\mathcal T})=\frac{x^2}{\tau}.
 \label{eq:cond-min-degenerate}
\end{equation}
For $m=1$, let $\lambda_{\rm diag}$ be the least curvature on the diagonal
part of $\mathcal T$ (directions diagonal in the cost eigenbases).  Then
\begin{equation}
 \frac{(dG-1)x^2+b(x)^2}{dG\tau}
 \le\lambda_{\rm diag}\le
 \frac{r(G-1)x^2+(x+\delta_{\min})^2}
      {\tau\{r(G-1)+1\}}.                                  \label{eq:cond-diagonal-sandwich}
\end{equation}
If $m=1<d$, the ground--excited winner direction, which couples the ground
eigenvector with an eigenvector of gap $\gamma_{\min}$, has curvature
$x(x+\gamma_{\min})/\tau$.  Consequently, with constants depending only on
the fixed spectral data,
\begin{equation}
 \lambda_{\min}(\mathcal Q|_{\mathcal T})\asymp\frac1\tau
 \begin{cases}
 \min\{x(x+\gamma_{\min}),
        x^2+(x+\delta_{\min})^2/G\},&m=1<d,\\
 x^2+(x+\delta_1)^2/G,&m=d=1.
 \end{cases}                                                \label{eq:cond-min-simple}
\end{equation}
For $G\ge3$,
\begin{equation}
 \frac{(x+\delta_{\max})^2}{\tau}
 \le\lambda_{\max}(\mathcal Q|_{\mathcal T})
 \le\frac{(x+\Lambda)^2}{\tau}.                           \label{eq:cond-max-sandwich}
\end{equation}
\end{lemma}

\begin{proof}
Use eigenbases of the two costs.  Diagonal and symmetric off-diagonal
coordinates are orthogonal invariant subspaces.  An off-diagonal $(i,j)$
coordinate in a block with shifted eigenvalues $d_i,d_j$ has curvature
$d_id_j/\tau$, and a diagonal coordinate has curvature $d_i^2/\tau$.
Only diagonal coordinates are coupled by \eqref{eq:cond-tangent}.

If $m\ge2$, the difference of two winner ground diagonals, or their
off-diagonal coordinate, is tangent and has curvature $x^2/\tau$.  No
ambient coordinate has smaller curvature, proving \eqref{eq:cond-min-degenerate}.
For $m=1$, write the winner's ground diagonal coefficient as $a$ and all other
$dG-1$ diagonal coefficients as $z$.  The trace constraint gives
$\mathbf1^\top z=-a$, hence $\norm{z}_2^2\ge a^2/(dG-1)$; every corresponding
curvature is at least $b(x)^2/\tau$.  Minimizing the resulting two-term quotient
gives the lower bound in \eqref{eq:cond-diagonal-sandwich}.  For the upper
bound, cancel $a$ equally among the $r(G-1)$ loser coordinates with gap
$\delta_{\min}$ and calculate its Rayleigh quotient.

The stated winner off-diagonal coordinate supplies the other upper bound in
\eqref{eq:cond-min-simple}; all remaining off-diagonal curvatures are bounded
below by the same expression up to fixed spectral constants.  Together with
the diagonal sandwich this proves \eqref{eq:cond-min-simple}, including the
scalar case where no within-block off-diagonal exists.  Finally, opposite
maximum-gap diagonal directions in two loser blocks are tangent and give the
lower bound in \eqref{eq:cond-max-sandwich}; its upper bound is the largest
curvature of an ambient coordinate.
\end{proof}

\begin{theorem}[Multiplicity-dependent conditioning phases]
\label{thm:cond-conditioning-phases}
For fixed $\alpha>0$, $\tau=\alpha/G$, and fixed spectral data,
\begin{equation}
 \kappa_{\rm red}=\begin{cases}
 \Theta(G^2),&m\ge2,\ 0<\alpha<\alpha_*,\\
 \Theta(G),&m=1,\ 0<\alpha\le\alpha_*,\\
 \Theta(G),&m\ge2,\ \alpha=\alpha_*,\\
 \Theta(1),&\alpha>\alpha_*\quad\text{for every }m.
 \end{cases}
 \label{eq:cond-kappa-phases}
\end{equation}
\end{theorem}

\begin{proof}
In the subcritical phase, \cref{thm:condensation-threshold} gives
$x=\Theta(G^{-1})$.  By \eqref{eq:cond-max-sandwich}, the largest
curvature is $\Theta(G)$.  For $m\ge2$,
\eqref{eq:cond-min-degenerate} gives minimum curvature $\Theta(G^{-1})$;
for $m=1$, \eqref{eq:cond-min-simple} gives $\Theta(1)$.  At criticality
$x=\Theta(G^{-1/2})$, and either minimum formula gives $\Theta(1)$, while
the maximum remains $\Theta(G)$.  Above capacity $x$ is bounded away from
zero, so every curvature is $\Theta(G)$.
\end{proof}

The ground--excited direction needs $m<d$.  For $m=d=1$, the diagonal
direction that spreads the winner's trace across the losers plays its role,
and \eqref{eq:cond-min-simple} gives the same orders of $\kappa_{\rm red}$.

\subsection{Finite mass--conditioning laws}
\label{sec:condensation-mass-conditioning}

The next theorem bounds $\kappa_{\rm red}$ at finite $G$ in terms of the
winner mass $p$, without fixing the schedule $\tau=\alpha/G$ of
\cref{thm:cond-conditioning-phases}.

\begin{theorem}[Finite mass--conditioning law]
\label{thm:cond-mass-conditioning}
Let $G\ge3$ and let $p=p_G\in(0,1)$ be the winner mass at an exact center.
Define
\begin{equation*}
 R_G(p)=\frac{(G-1)p}{1-p}. \condtag{MC1}\label{eq:cond-mc1}
\end{equation*}
Then
\begin{align*}
 m\ge2:&\qquad \kappa_{\rm red}\ge R_G(p)^2,              \condtag{MC2}\label{eq:cond-mc2}\\
 m=1<d:&\qquad \kappa_{\rm red}\ge
 \min\left\{1,\frac{\delta_{\max}}{\gamma_{\min}}\right\}R_G(p).
                                                               \condtag{MC3}\label{eq:cond-mc3}
\end{align*}
If $m=d=1$, writing $\delta=\delta_1$, the exact closed form is
\begin{equation*}
 \boxed{\displaystyle
 \kappa_{\rm red}=\frac{G}
 {1+(1-p)^2/\{(G-1)p^2\}}.}                                \condtag{MC4}\label{eq:cond-mc4}
\end{equation*}
For scalar blocks with $G=2$, the tangent space is one-dimensional and
$\kappa_{\rm red}=1$, so \eqref{eq:cond-mc4} is asserted only for $G\ge3$.
\end{theorem}

\begin{proof}
The exact mass equations are
\begin{equation}
 p=\tau A(x),\qquad 1-p=\tau(G-1)B(x).                    \label{eq:cond-mass-equations}
\end{equation}
Because $A(x)\le d/x$ and $B(x)\ge d/(x+\delta_{\max})$,
\begin{equation}
 \frac{1-p}{p}=(G-1)\frac{B(x)}{A(x)}
 \ge(G-1)\frac{x}{x+\delta_{\max}},
 \qquad \frac{x+\delta_{\max}}x\ge R_G(p).              \label{eq:cond-mass-ratio}
\end{equation}
The tangent direction with opposite unit entries on the $\delta_{\max}$
diagonal coordinates of two loser blocks gives
$\lambda_{\max}\ge(x+\delta_{\max})^2/\tau$.  If $m\ge2$, divide this
by \eqref{eq:cond-min-degenerate} and use
\eqref{eq:cond-mass-ratio}.

If $m=1<d$, the winner ground--excited direction gives
$\lambda_{\min}\le x(x+\gamma_{\min})/\tau$, while
\[
 \frac{x+\delta_{\max}}{x+\gamma_{\min}}
 \ge\min\left\{1,\frac{\delta_{\max}}{\gamma_{\min}}\right\}.
\]
Factoring the curvature ratio proves \eqref{eq:cond-mc3}.

For $d=m=1$, put $a=x^2/\tau$ and $b=(x+\delta)^2/\tau$.  On the scalar
trace-zero space, $G-2$ loser-difference directions have curvature $b$, and
the direction balancing the winner against equal loser changes has curvature
$\varrho=\{(G-1)a+b\}/G$.  Since $b>a$, these are the largest and smallest
eigenvalues.  The mass equations reduce to
\[
 \frac{x}{x+\delta}=\frac{1-p}{(G-1)p}.
\]
Thus $\kappa_{\rm red}=b/\varrho=G/\{1+(G-1)a/b\}$, which is
\eqref{eq:cond-mc4}.
\end{proof}

In particular, $p\ge p_0>0$ forces $\kappa_{\rm red}=\Omega(G^2)$ when
$m\ge2$ and $\kappa_{\rm red}=\Omega(G)$ when $m=1$.  This also follows from
$x\le d\tau/p_0$, the bound $\tau=O(G^{-1})$ from the loser mass, and the
test directions above; the exact laws add the explicit dependence on $p$.

\subsection{Visibility, collision, and the conditioning dichotomy}
\label{sec:condensation-visibility}

The next result uses only the symmetry of the two-cost model, not the
spectral assumptions of \cref{sec:condensation-model}.  Let
$C_W,C_L\in\Sym^d$, $G\ge2$, and $\tau>0$.  We compare the trace-one center
$\rho_0$ where every block has the loser cost with the center $\rho_1$ where
block $1$ has cost $C_W$:
\begin{align}
 \rho_0&=\arg\min_{\substack{Z_g\succ0\\\sum_g\tr Z_g=1}}
 \left\{\sum_{g=1}^G\inner{C_L}{Z_g}-\tau\sum_g\log\det Z_g\right\},
                                                               \label{eq:cond-rho-zero}\\
 \rho_1&=\arg\min_{\substack{Z_g\succ0\\\sum_g\tr Z_g=1}}
 \left\{\inner{C_W}{Z_1}+\sum_{g=2}^G\inner{C_L}{Z_g}
       -\tau\sum_g\log\det Z_g\right\}.                 \label{eq:cond-rho-one}
\end{align}
Identify each tuple with its block-diagonal density matrix and write
\begin{equation}
 \rho_0=\bigoplus_{g=1}^G L_0,\qquad
 \rho_1=W_1\oplus\bigoplus_{g=2}^G L_1,\qquad p=\tr W_1.
 \label{eq:cond-rho-blocks}
\end{equation}
Permutation symmetry gives $\tr L_0=1/G$ and
$\tr L_1=(1-p)/(G-1)$.

\begin{theorem}[Visibility and label collision]
\label{thm:cond-visibility}
For the centers above,
\begin{equation*}
 \boxed{\left|p-\frac1G\right|
 \le\Dtr(\rho_0,\rho_1)\le\max\left\{p,\frac1G\right\}.}
 \condtag{O1}\label{eq:cond-visibility}
\end{equation*}
If the block labels of two independent copies are measured and compared,
the probabilities $c_0,c_1$ that they agree are
\begin{equation*}
 c_0=\frac1G,\qquad
 c_1=p^2+\frac{(1-p)^2}{G-1},\qquad
 \boxed{c_1-c_0=\frac{(Gp-1)^2}{G(G-1)}.}
 \condtag{O2}\label{eq:cond-collision}
\end{equation*}
The measurement is public (independent of the input) and acts only on the
two label registers, so the identity holds regardless of degeneracy inside
the blocks and is unchanged by purifying and then discarding auxiliary
registers.
\end{theorem}

\begin{proof}
Measuring whether the label is $1$ gives the lower bound in
\eqref{eq:cond-visibility}.  Let $\lambda_0,\lambda_1$ be the scalar trace
multipliers of the two centers.  Their loser blocks satisfy
\[
 L_b=\tau(C_L-\lambda_bI)^{-1}\quad(b=0,1),
\]
and are therefore Loewner ordered.  If $\lambda_1\le\lambda_0$, then
$L_1\preceq L_0$ and the trace identities imply $p\ge1/G$.  Hence
\[
 (G-1)\norm{L_1-L_0}_1=p-\frac1G,
 \qquad \norm{W_1-L_0}_1\le p+\frac1G.
\]
Adding the trace norms of the direct-sum blocks and dividing by two gives
$\Dtr(\rho_0,\rho_1)\le p$.  If $\lambda_1\ge\lambda_0$, the order and
inequality reverse, and the same calculation gives the upper bound $1/G$.
Finally the two label distributions are
\[
 (1/G,\ldots,1/G),\qquad
 \left(p,\frac{1-p}{G-1},\ldots,\frac{1-p}{G-1}\right).
\]
Summing their squared entries and subtracting proves
\eqref{eq:cond-collision}.
\end{proof}

Only the upper trace-distance bound uses the resolvent form of the centers;
the lower bound and the collision identity need only a uniform label
distribution for $\rho_0$ and exchangeable losers.  The condition $p\ge1/G$
holds whenever $C_W\preceq C_L$, since $X\mapsto(X-\lambda I)^{-1}$ is
operator antitone on $X\succ\lambda I$, and for the family of
\cref{sec:wta-central-path} by \eqref{eq:wta-resolvent-order}.  Then the
trace-distance bounds reduce to $p-1/G\le\Dtr\le p$.

For $C_W=C_*$, under the spectral assumptions of
\cref{sec:condensation-model} and $G\ge3$,
\cref{thm:cond-visibility,thm:cond-mass-conditioning} give the following
finite laws relating conditioning and visibility.  With
$c=\min\{1,\delta_{\max}/\gamma_{\min}\}$,
\begin{align*}
 p&\le\frac{\sqrt{\kappa_{\rm red}}}
 {G-1+\sqrt{\kappa_{\rm red}}} &&(m\ge2),                 \condtag{CV1}\label{eq:cond-cv1}\\
 p&\le\frac{\kappa_{\rm red}}
 {c(G-1)+\kappa_{\rm red}} &&(m=1<d),                    \condtag{CV2}\label{eq:cond-cv2}\\
 p&=\frac1{1+\sqrt{G-1}\sqrt{G/\kappa_{\rm red}-1}}
 &&(m=d=1).                                                \condtag{CV3}\label{eq:cond-cv3}
\end{align*}
By \eqref{eq:cond-visibility}, the maximum of each right-hand side and $1/G$
bounds $\Dtr(\rho_0,\rho_1)$.  Conversely, if
$\Dtr(\rho_0,\rho_1)\ge\upsilon>1/G$, then $p\ge\upsilon$ and
\begin{equation*}
 \kappa_{\rm red}\ge
 \left(\frac{(G-1)\upsilon}{1-\upsilon}\right)^2\ (m\ge2),
 \qquad
 \kappa_{\rm red}\ge
 c\frac{(G-1)\upsilon}{1-\upsilon}\ (m=1<d).             \condtag{CV4}\label{eq:cond-cv4}
\end{equation*}
For $m=d=1$, \eqref{eq:cond-cv3} gives the same linear order.  Thus
constant visibility forces $\kappa_{\rm red}=\Omega(G^2)$ when $m\ge2$ and
$\kappa_{\rm red}=\Omega(G)$ when $m=1$.  A constant collision gap also
forces $p=\Omega(1)$ and hence the same dichotomy.

\paragraph{Approximate outputs and reduced symmetry.}
Let $\widetilde\rho_b$ satisfy
$\Dtr(\widetilde\rho_b,\rho_b)\le\epsilon$.  Contractivity and the
tensor-product triangle inequality give
$\Dtr(\widetilde\rho_b^{\otimes2},\rho_b^{\otimes2})\le2\epsilon$.
Thus the observed two-copy collision gap is at least
\begin{equation}
 \Delta_{\rm out}\ge d(p,G)-4\epsilon,
 \qquad d(p,G):=\frac{(Gp-1)^2}{G(G-1)}.                  \label{eq:cond-output-gap}
\end{equation}
By standard concentration bounds, $O(\Delta_{\rm out}^{-2})$ repetitions
distinguish the two cases with constant success probability, and an extra
$\log(1/\delta)$ factor gives failure probability $\delta$.

The collision lower bound does not need equal loser masses.  Suppose a
public label measurement is uniform on the no-winner input and, on every
unique-winner input, gives label masses $w_g$ with winner mass
$q\ge p\ge1/G$.  Cauchy--Schwarz gives
\begin{equation}
 \sum_{g=1}^G w_g^2\ge q^2+\frac{(1-q)^2}{G-1}
 \ge p^2+\frac{(1-p)^2}{G-1},                            \label{eq:cond-symmetry-light}
\end{equation}
where the last inequality holds because
$q\mapsto q^2+(1-q)^2/(G-1)$ is increasing for $q\ge1/G$.  The collision gap is therefore at least $d(p,G)$, with
equality exactly at winner mass $p$ and uniform loser mass.  If the
no-winner collision probability is only bounded above by $b$, replace the
gap by $p^2+(1-p)^2/(G-1)-b$.  In particular, if $p\ge p_0>0$ and
$G\ge2/p_0$, then $d(p,G)\ge p_0^2/4$.

\subsection{Unique-OR query lower bounds}
\label{sec:condensation-query-lower}

To turn \cref{thm:cond-visibility} into a query lower bound, we use a
composed problem.  Let $G\ge2$ and let a possibly partial Boolean predicate
$f:\mathcal X_0\sqcup\mathcal X_1\to\{0,1\}$ have general adversary value
$\mathsf{Adv}_f=\operatorname{Adv}^{\pm}(f)$.  For $G$ independently queried
blocks, set
\begin{equation}
 \mathcal P_0=\mathcal X_0^G,
 \qquad
 \mathcal P_1=\bigsqcup_{a=1}^G
 \mathcal X_0^{a-1}\times\mathcal X_1\times\mathcal X_0^{G-a},
 \label{eq:cond-unique-or-promise}
\end{equation}
and let $F$ be the problem of distinguishing $\mathcal P_0$ (no winner) from
$\mathcal P_1$ (exactly one winner).

\begin{theorem}[Unique-OR adversary and preparation]
\label{thm:cond-unique-or}
In the raw-query model where one query addresses one coordinate of one block,
\begin{equation*}
 \operatorname{Adv}^{\pm}(F)\ge \mathsf{Adv}_f\sqrt G.     \condtag{O3}\label{eq:cond-unique-or-adv}
\end{equation*}
Consequently, suppose that fixed public projectors onto the blocks give
the uniform label distribution on every no-winner target state and winner
mass at least $p\ge1/G$ on every unique-winner target state; loser masses
may be arbitrary.  Any circuit that prepares, on every promised input, a state
within trace distance $\epsilon$ of such a target uses
$\Omega(\mathsf{Adv}_f\sqrt G)$ raw queries whenever
$d(p,G)-4\epsilon=\Omega(1)$.  By tightness of the general adversary bound,
this is equivalently $\Omega(Q(f)\sqrt G)$ up to universal constants.
\end{theorem}

\begin{proof}
Take the rectangular zero-to-one block $B$, with rows indexed by
$\mathcal X_0$ and columns by $\mathcal X_1$, of a feasible adversary matrix
for $f$, normalized so that
\begin{equation}
 \norm{B}=\mathsf{Adv}_f,\qquad
 \max_j\norm{B\circ\Delta_j}\le1,                         \label{eq:cond-inner-adversary}
\end{equation}
where $\Delta_j[x,y]=1$ exactly when query $j$ has different answers on
$x,y$, and let $I_0$ be the identity on $\ell_2(\mathcal X_0)$.  For the
sector of $\mathcal P_1$ whose winner is block $a$, define
\[
 C_a=I_0^{\otimes(a-1)}\otimes B\otimes I_0^{\otimes(G-a)},
 \qquad C=[C_1\ C_2\ \cdots\ C_G],
\]
and use the Hermitian adversary matrix
\[
 \Gamma=\begin{pmatrix}0&C\\C^*&0\end{pmatrix}.
\]
The winner sectors are disjoint, and
\[
 CC^*=\sum_{a=1}^G I_0^{\otimes(a-1)}\otimes BB^*\otimes
                          I_0^{\otimes(G-a)}.
\]
The summands commute.  Tensoring a top eigenvector of $BB^*$ across all
$G$ blocks gives their common top eigenvector, so
$\norm{\Gamma}=\mathsf{Adv}_f\sqrt G$.  Filtering by raw query $(a,j)$ kills every
winner sector except $a$, since row and column inputs agree in block $a$ in
all other sectors.  The remaining block is $B\circ\Delta_j$ and has norm at
most one by \eqref{eq:cond-inner-adversary}.  Thus $\Gamma$ is feasible and
proves \eqref{eq:cond-unique-or-adv}.

If $d(p,G)-4\epsilon$ is a positive constant,
\eqref{eq:cond-output-gap} turns a constant number of fresh preparations and
public collision measurements into a bounded-error algorithm for $F$.
The adversary lower bound therefore applies to the preparation circuit.
Tightness and composition of the negative-weight adversary bound give the
equivalent $Q(f)$ statement
\cite{hyer2005-tight-adversary-bounds-for-composite,lee2011-quantum-query-complexity-of-state,reichardt2011-reflections-for-quantum-query-algorithms}.
\end{proof}

The theorem assumes that all product inputs in
\eqref{eq:cond-unique-or-promise} are valid, both fibers of $f$ are
nonempty, and the block projectors are public.  It needs no clean evaluator, candidate
verification, inverse preparation, or restriction on the internal states of
the blocks.  If one coefficient-oracle call can be
simulated with at most $s$ raw queries, the theorem gives a lower bound of
\begin{equation}
 \Omega\!\left(\frac{\mathsf{Adv}_f\sqrt G}{s}\right)     \label{eq:cond-coefficient-transfer}
\end{equation}
coefficient-oracle calls.  An oracle that returns a winner bit, an
aggregate flag, or an input-dependent projector is stronger and may remove
the inner factor $\mathsf{Adv}_f$.

No such lower bound holds when the output error is large compared with the
visibility: by \cref{thm:cond-visibility}, outputting the public all-loser
center $\rho_0$ uses no queries and is valid on every input with zero or
one winner when
\begin{equation}
 \epsilon\ge\max\{p,1/G\}.                                \label{eq:cond-zero-query-boundary}
\end{equation}
Hence a constant primal objective gap \eqref{eq:objective-gap-g} need not imply a lower bound for
trace-normalized primal-density output at constant error.  This concerns
the direct-sum density of \cref{def:trace-density-output}, not a normalized
vectorization, whose label probabilities are Frobenius rather than trace
weights.

\subsection{Preparation, extraction, and the conservation law}
\label{sec:condensation-algorithms}

For the upper bounds we assume exactly one winner $w$ and a clean coherent
evaluator of the predicate
\begin{equation*}
 V_f\ket{g,0}=\ket{g,f_g},\qquad f_g=1\Longleftrightarrow g=w,          \condtag{A1}\label{eq:cond-clean-evaluator}
\end{equation*}
that costs $T_{\rm eval}$ raw queries, as does its inverse.  The two cost
matrices, $G$, and $\tau$ are public, so the center has the form
\begin{equation*}
 \rho_p=p\ket w\!\bra w\otimes\rho_W
 +\frac{1-p}{G-1}\sum_{g\ne w}\ket g\!\bra g\otimes\rho_L,            \condtag{A2}\label{eq:cond-central-state}
\end{equation*}
where the normalized block resolvents $\rho_W,\rho_L$ are public.  Put
\begin{equation*}
 a_W=p,\qquad a_L=\frac{1-p}{G-1},\qquad
 a_{\max}=\max\{a_W,a_L\}.                                            \condtag{A3}\label{eq:cond-component-weights}
\end{equation*}

\begin{proposition}[Coherent rejection preparation]
\label{prop:cond-rejection-preparation}
There is a purification circuit for \eqref{eq:cond-central-state} using
\begin{equation*}
 \Otilde\!\left(T_{\rm eval}\sqrt{Ga_{\max}}
              +T_{\rm loc}\right)                                    \condtag{A4}\label{eq:cond-preparation-upper}
\end{equation*}
operations, where $T_{\rm loc}$ is the gate cost (not a raw-query cost) of
the one public preparation of $\rho_W$ or $\rho_L$ conditioned on the
predicate.  When $p\ge1/G$, the raw-query part is
\begin{equation*}
 \Otilde(T_{\rm eval}\sqrt{Gp}).                                     \condtag{A5}\label{eq:cond-preparation-upper-p}
\end{equation*}
\end{proposition}

\begin{proof}
Prepare $G^{-1/2}\sum_g\ket g$, compute $f_g$, and rotate an acceptance
qubit with amplitude $\sqrt{a_{f_g}/a_{\max}}$, but do not yet copy the
label or prepare the internal state.  This trial succeeds with probability
\begin{equation*}
 P_{\rm acc}=\frac1G\left(\frac{a_W}{a_{\max}}+
                    (G-1)\frac{a_L}{a_{\max}}\right)
             =\frac1{Ga_{\max}}.                                      \condtag{A6}\label{eq:cond-acceptance}
\end{equation*}
This probability is known because the secular equation is solved
classically, so exact amplitude amplification with matched phases applies
to this trial, acting on the label, predicate, and acceptance registers
\cite{brassard2002-quantum-amplitude-amplification-and-estimation}.  After
it succeeds, the label has probability $a_{f_g}$; only then coherently copy
the label to an environment register and use the retained predicate once to
prepare a purification of $\rho_W$ or $\rho_L$.  Tracing out the copied
label, predicate, acceptance, and local purification registers gives
exactly \eqref{eq:cond-central-state}.  Only the evaluator calls are
repeated by the amplification, and the local preparation runs once, proving
both bounds; the tilde hides only polylogarithmic factors in the requested
implementation accuracy.
\end{proof}

The secular equation and the fixed-size resolvents are computed
classically; if $d$ varies, one must add $\poly(d,\log(1/\epsilon))$ costs
for arithmetic and state synthesis.  A bounded-error evaluator must first
have its error reduced coherently below the final state error.

Suppose independent copies, each costing $Q_\rho$ raw queries, satisfy
$\Dtr(\widetilde\rho,\rho_p)\le\epsilon<p$.  Measuring the label proposes
the winner with probability at least $p-\epsilon$, so verifying and
repeating finds the winner with failure probability $\delta$ using
\begin{equation*}
 O\!\left(\frac{Q_\rho+T_{\rm eval}}{p-\epsilon}
          \log\frac1\delta\right)                                     \condtag{A7}\label{eq:cond-copy-extraction}
\end{equation*}
raw queries.  For exact copies take $\epsilon=0$; if
$\epsilon\le p/2$, this has the usual $O(1/p)$ repetition factor.  Given
instead a coherent preparation unitary $U_\rho$ costing $Q_U$ raw queries,
its inverse, and a clean winner reflection, amplitude amplification gives
\begin{equation*}
 O\!\left(\frac{Q_U+T_{\rm eval}}{\sqrt p}\log\frac1\delta\right).     \condtag{A8}\label{eq:cond-coherent-extraction}
\end{equation*}
This bound needs $U_\rho$ and its inverse as operators; an accurate output
density matrix says nothing about how the unitary acts on inputs other than
$\ket0$.

Combining \eqref{eq:cond-preparation-upper-p} and
\eqref{eq:cond-coherent-extraction} gives the preparation--extraction
conservation law
\begin{equation*}
 \Otilde(T_{\rm eval}\sqrt{Gp})\,O(p^{-1/2})
   =\Otilde(T_{\rm eval}\sqrt G).                                     \condtag{A9}\label{eq:cond-conservation}
\end{equation*}
Stronger condensation moves work from extraction to preparation without
changing the end-to-end number of raw queries up to logarithmic factors.
The three regimes give the following costs for exact output
($\epsilon\le p/2$ suffices in the copy-extraction column):
\begin{equation*}
\begin{array}{@{}lccc@{}}
 \text{regime} & \text{central preparation} & \text{copy extraction}
   & \text{coherent extraction}\\ \hline
 \text{subcritical} & \Otilde(T_{\rm eval}G^{1/2}) & O(1) & O(1)\\
 \text{critical} & \Otilde(T_{\rm eval}G^{1/4}) & O(G^{1/2}) & O(G^{1/4})\\
 \text{supercritical} & \Otilde(T_{\rm eval}) & O(G) & O(G^{1/2})
\end{array}                                                            \condtag{A10}\label{eq:cond-algorithm-phases}
\end{equation*}
The supercritical preparation bound matters only for exact output or for
accuracy fine enough to resolve a winner mass of order $1/G$; at fixed
error, the zero-query output of \eqref{eq:cond-zero-query-boundary} may
suffice.

Alternatively, quantum search finds the winner at cost
\begin{equation*}
 \Otilde(T_{\rm eval}\sqrt G),                                      \condtag{A11}\label{eq:cond-direct-search}
\end{equation*}
and solving and recovering the selected block at costs $T_{\rm solve}$ and
$T_{\rm recover}$ gives the end-to-end bound
\begin{equation*}
 \Otilde(T_{\rm eval}\sqrt G)+T_{\rm solve}+T_{\rm recover}.          \condtag{A12}\label{eq:cond-search-crossover}
\end{equation*}
These bounds use standard quantum search and minimum finding
\cite{boyer1998-tight-bounds-on-quantum-searching,durr1996-a-quantum-algorithm-for-finding}.
This is search followed by crossover, not a new interior-point method.  If
the local predicate has bounded-error raw-query complexity $q$ and the
blocks are independent, composition of the general adversary bound
\cite{hyer2005-tight-adversary-bounds-for-composite,lee2011-quantum-query-complexity-of-state,reichardt2011-reflections-for-quantum-query-algorithms}
gives the total bounded-error raw-query cost
\begin{equation*}
 Q(\operatorname{UniqueOR}_G\circ f)=\Theta(q\sqrt G).                \condtag{A13}\label{eq:cond-search-composition}
\end{equation*}
Combining this lower bound with the coherent extraction bound
\eqref{eq:cond-coherent-extraction} gives
\begin{equation*}
 Q_U+T_{\rm eval}=\Omega(q\sqrt{Gp}),                                  \condtag{A14}\label{eq:cond-prep-tradeoff}
\end{equation*}
and, for $\epsilon$-accurate independent copies with $\epsilon<p$, the copy
extraction bound gives the weaker
\begin{equation*}
 Q_\rho+T_{\rm eval}=\Omega((p-\epsilon)q\sqrt G).                    \condtag{A15}\label{eq:cond-copy-tradeoff}
\end{equation*}
For exact output this is $\Omega(pq\sqrt G)$, and the same asymptotic bound
holds when $\epsilon\le p/2$.
These statements require independent block oracles; if one query can touch
several blocks, or an aggregate winner flag is available, the composition
argument fails.

For comparison, the block product barrier has parameter $dG$, so short-step
path following from constant $\tau$ to $\tau=\Theta(1/G)$ takes
$O(\sqrt{dG}\log G)$ Newton iterations.  This count excludes the linear
solves, whose cost depends on a condition number, and state preparation and
recovery.  It must not be multiplied by \eqref{eq:cond-kappa-phases} unless
an actual solver and preconditioner are chosen and analyzed.

\subsection{Sublinear multiwinner capacity}
\label{sec:condensation-multiwinner}

Now let exactly $k$ blocks have cost $C_W$, with
$v=\lambda_{\min}(C_W)$ of multiplicity $m$, and suppose
$C_L-vI\succ0$.  In this subsection, redefine
\begin{equation}
 A(x)=\tr(C_W-vI+xI)^{-1},\qquad
 B(x)=\tr(C_L-vI+xI)^{-1}.                              \label{eq:cond-multi-resolvents}
\end{equation}
Then $A(x)=m/x+O(1)$,
$B_0=B(0)<\infty$, and $\beta=-B'(0)>0$.  The exact equations are
\begin{equation}
 1=\tau\{kA(x)+(G-k)B(x)\},\qquad
 p_{G,k}=\tau kA(x),\quad1-p_{G,k}=\tau(G-k)B(x).         \label{eq:cond-multi-secular}
\end{equation}

\begin{theorem}[Multiwinner capacity law]
\label{thm:cond-multiwinner}
Let $1\le k=k_G=o(G)$, put $\tau=\alpha/G$, and keep $\alpha>0$ fixed.
With $\alpha_*=1/B_0$,
\begin{align}
 p_{G,k}&\longrightarrow1-\alpha B_0,
 &x&\sim\frac{\alpha m k}{G(1-\alpha B_0)}
 &&(0<\alpha<\alpha_*),                                  \label{eq:cond-multi-subcritical}\\
 p_{G,k}&\sim\alpha_*\sqrt{m\beta k/G},
 &x&\sim\sqrt{mk/(\beta G)}
 &&(\alpha=\alpha_*),                                    \label{eq:cond-multi-critical}\\
 p_{G,k}&\sim\alpha A(x_\alpha)k/G,
 &B(x_\alpha)&=1/\alpha
 &&(\alpha>\alpha_*).                                    \label{eq:cond-multi-supercritical}
\end{align}
\end{theorem}

\begin{proof}
Substitute $\tau=\alpha/G$ into the first equation in
\eqref{eq:cond-multi-secular}.  Below criticality $x\to0$: otherwise the
winner term vanishes because $k/G\to0$, while the loser term is at most
$\alpha B_0<1$.  The loser mass then tends to $\alpha B_0$, proving the
first mass law, and the pole $A(x)\sim m/x$ gives the formula for $x$.

At criticality put $q=k/G$.  Taylor expansion of the loser equation and the
winner pole give two expressions for the total winner mass:
\begin{equation}
 p_{G,k}=q+\frac{\beta}{B_0}x+o(q+x)
       =\frac{qm}{B_0x}+O(q).                            \label{eq:cond-multi-balance}
\end{equation}
As in the one-winner proof, monotonicity first bounds $x$ by constant
multiples of $\sqrt q$; balancing the leading terms then gives
$x\sim\sqrt{mq/\beta}$ and the critical mass.  Above criticality the same
compactness and strict-monotonicity argument as in
\cref{thm:condensation-threshold} gives $x\to x_\alpha>0$; the winner
equation gives the final law.
\end{proof}

The assumption $k=o(G)$ is needed.  If $k/G\to q\in(0,1]$, the limit
instead solves $1=\alpha\{qA(x)+(1-q)B(x)\}$ with $x>0$, and there is no
singular capacity transition at fixed $\alpha$.

\paragraph{Nonidentical losers.}
Let every losing cost satisfy $C_{\ell,G}-vI\succ0$, and define
\[
 B_{\ell,G}(x)=\tr(C_{\ell,G}-vI+xI)^{-1}
\]
and their average
\[
 \overline B_G(x)=\frac1{G-k}
       \sum_{\ell\ {\rm losing}}B_{\ell,G}(x).
\]
If $k=o(G)$ and $\overline B_G\to B$ locally uniformly on $x\ge0$, with
$B(0)=B_0\in(0,\infty)$, the subcritical and supercritical conclusions and
the subcritical formula for $x$ remain valid.  The critical law also holds
if, in addition,
\begin{equation}
 \overline B_G(0)-B_0=o(\sqrt{k/G}),\qquad
 \overline B_G(x)=\overline B_G(0)-\beta x+r_G(x),         \label{eq:cond-empirical-resolvent}
\end{equation}
where $\beta>0$ and, for every fixed $C>0$,
\[
 \sup_{0<x\le C\sqrt{k/G}}\frac{|r_G(x)|}{x}\longrightarrow0.
\]
Indeed, replacing $B$ by $\overline B_G$ in
\eqref{eq:cond-multi-secular} leaves the subcritical and supercritical limit
arguments unchanged.  Under \eqref{eq:cond-empirical-resolvent}, the
finite-size error is of lower order than both leading terms in
\eqref{eq:cond-multi-balance}, which proves the critical law.  These
hypotheses control only the winner mass; they do not give exchangeable
losers, a common public internal state, or the exact collision identity.

For identical blocks and a fixed winner set $S$ of size $k<G$, let
$q=k/G$.  The Loewner-order argument in the proof of
\cref{thm:cond-visibility} gives
\begin{equation}
 |p_{G,k}-q|\le\Dtr(\rho_0,\rho_S)\le\max\{p_{G,k},q\},     \label{eq:cond-multi-visibility}
\end{equation}
and direct calculation gives the two-copy label-collision gap
\begin{equation}
 c_S-c_0=\frac{p_{G,k}^2}{k}+\frac{(1-p_{G,k})^2}{G-k}-\frac1G
 =\boxed{\frac{(Gp_{G,k}-k)^2}{Gk(G-k)}}.                 \label{eq:cond-multi-collision}
\end{equation}
Thus a constant total winner mass gives a collision gap of only $O(1/k)$
when $k$ grows, and a trace-distance bound for each fixed $S$ does not give
one measurement for an unknown $S$.  \Cref{sec:state-conversion}
avoids this $1/k$ loss for its structured family by the composed
small-success bound (\cref{thm:sc-small-success}); that stronger conclusion
is not available in the abstract two-cost model.

\subsection{Exactly-\texorpdfstring{$k$}{k} coherent preparation}
\label{sec:condensation-exact-k-preparation}

Assume now that $1\le k<G$ is known, that there are exactly $k$ winners, and
that a clean coherent membership marker costs $T_{\rm mark}$ raw queries.
The two costs, $G$, $k$, $\tau$, the trace multiplier, and the normalized
internal resolvents are public, and controlled marker calls and the inverse
preparation circuit are available.  Put, for $p=p_{G,k}$,
\begin{equation*}
 a_W=\frac pk,\qquad a_L=\frac{1-p}{G-k},\qquad
 a_{\max}=\max\{a_W,a_L\}.                               \condtag{AP0}\label{eq:cond-ap0}
\end{equation*}
The construction in \cref{prop:cond-rejection-preparation}, with the
membership bit in place of the unique-winner bit, succeeds in one trial with
probability $1/(Ga_{\max})$.  This probability is known exactly, so exact
amplitude amplification with matched phases uses
\begin{equation*}
 O(\sqrt{Ga_{\max}}+1)                                    \condtag{AP1}\label{eq:cond-ap1}
\end{equation*}
marker calls
\cite{brassard2002-quantum-amplitude-amplification-and-estimation}.  After
it succeeds, we copy the label to an environment register and, controlled
on the membership bit, prepare once a public purification of the normalized
winner or loser block, where $W$ and $L$ are the common unnormalized winner
and loser blocks of the center and
\begin{equation*}
 \omega_W=\frac{W}{p/k},\qquad
 \omega_L=\frac{L}{(1-p)/(G-k)}.                                     \condtag{AP2}\label{eq:cond-ap2}
\end{equation*}
The resulting circuit satisfies
\begin{equation*}
 \boxed{Q_U=O\!\left(T_{\rm mark}[\sqrt{Ga_{\max}}+1]\right).}         \condtag{AP3}\label{eq:cond-ap3}
\end{equation*}
When $p\ge k/G$, this is $O(T_{\rm mark}[\sqrt{Gp/k}+1])$; when
$p<k/G$, it is $O(T_{\rm mark}[\sqrt{G(1-p)/(G-k)}+1])$.  Preparing to
error $\epsilon>0$ rather than exactly adds the standard factor logarithmic
in $1/\epsilon$.  Costs of loading nonpublic eigenvectors, resolvents, or
multipliers must be added.

Conversely, amplification from the prepared purification finds a winner in
\begin{equation*}
 O\!\left(\frac{Q_U+T_{\rm mark}}{\sqrt p}\right)          \condtag{AP4}\label{eq:cond-ap4}
\end{equation*}
raw queries.  If the local predicate has bounded-error raw-query complexity
$q_f$ and independent composition gives the lower bound
$\Omega(q_f\sqrt{G/k})$ for finding a winner, then \eqref{eq:cond-ap4}
implies
\begin{equation*}
 \boxed{Q_U+T_{\rm mark}=\Omega\!\left(q_f\sqrt{Gp/k}\right).}         \condtag{AP5}\label{eq:cond-ap5}
\end{equation*}
This requires the preparation unitary and its inverse as operators; for
mixed-state output, independent copies give only the weaker tradeoff from
repetition.

If $f$ is the parity of $N$ input bits, then $T_{\rm mark}=\Theta(N)$ and
adversary composition gives winner-finding complexity
$\Theta(N\sqrt{G/k})$.  For $k=o(G)$, \cref{thm:cond-multiwinner} gives the
parity profile
\begin{equation*}
 Q_U=\begin{cases}
 \Theta(N\sqrt{G/k}),&0<\alpha<\alpha_*,\\
 \Theta(N(G/k)^{1/4}),&\alpha=\alpha_*,\\
 O(N),&\alpha>\alpha_*.
 \end{cases}                                                \condtag{AP6}\label{eq:cond-ap6}
\end{equation*}
The first two lines match \eqref{eq:cond-ap5} when the marker is optimal and
the amplification factor diverges.  The supercritical upper bound follows
from $a_{\max}=\Theta(1/G)$ even if $p<k/G$.  At criticality $p\to0$, so
this operator-level lower bound need not hold for density output at fixed
error.

This circuit does not solve the problem under the promise of zero or $k$
winners: with no winner, the center has a different trace multiplier and
loser resolvent.  Under that promise, exact existence search followed by
conditional preparation costs $O(T_{\rm mark}\sqrt{G/k})$, while returning
the public all-loser center costs zero queries and has error at most
$\max\{p_{G,k},k/G\}$ by \eqref{eq:cond-multi-visibility}.  The two-copy
collision measurement gives a matching lower bound only when the gap
\eqref{eq:cond-multi-collision}, reduced by the output error, stays
constant.

\subsection{Robust approximate centrality}
\label{sec:condensation-robustness}

The mass and conditioning bounds also hold, with modified constants, at
approximately central points whose blocks need not commute with the costs.
Keep the one-winner notation and, for positive definite $\widetilde Z_g$
and $\widetilde\lambda<v_*$, put
\begin{equation}
 x=v_*-\widetilde\lambda,\quad
 D_*=C_*-\widetilde\lambda I,\quad D_L=C_L-\widetilde\lambda I,\quad
 \widetilde T=\sum_g\tr\widetilde Z_g,\quad
 \widetilde p=\frac{\tr\widetilde Z_*}{\widetilde T}.       \label{eq:cond-robust-defs}
\end{equation}
Write $D_g=D_*$ for the winner and $D_g=D_L$ for every loser.
Assume a dimensionless stationarity residual in the barrier norm and a
trace residual:
\begin{equation}
 E_g=\tau^{-1}\widetilde Z_g^{1/2}D_g\widetilde Z_g^{1/2}-I,\qquad
 \max_g\norm{E_g}_{\rm op}\le\eta<1,\qquad
 |\widetilde T-1|\le\zeta<1.                              \label{eq:cond-robust-residual}
\end{equation}

\begin{theorem}[Robust capacity, mass, and conditioning]
\label{thm:cond-robust}
Every point satisfying \eqref{eq:cond-robust-defs}--
\eqref{eq:cond-robust-residual} satisfies the noncommutative Loewner
sandwich
\begin{equation*}
 (1-\eta)\tau D_g^{-1}\preceq\widetilde Z_g
 \preceq(1+\eta)\tau D_g^{-1}.                            \condtag{R1}\label{eq:cond-r3}
\end{equation*}
Consequently,
\begin{equation*}
 \boxed{\widetilde p\ge
 1-\frac{(1+\eta)\tau(G-1)B_0}{1-\zeta}
 \ge1-\frac{(1+\eta)\alpha B_0}{1-\zeta}.}               \condtag{R2}\label{eq:cond-r4}
\end{equation*}
Hence a uniform margin $(1+\eta)\alpha B_0\le1-\zeta-c$ with a constant
$c>0$ forces a positive, $G$-independent lower bound on the winner mass.  Conversely, if $G\to\infty$,
$\alpha\to\bar\alpha$, and $x\to0$, then
\begin{equation*}
 (1-\eta)\bar\alpha B_0\le1+\zeta.                       \condtag{R3}\label{eq:cond-r5}
\end{equation*}
Thus $(1+\eta)\alpha B_0=1-\zeta$ and
$(1-\eta)\alpha B_0=1+\zeta$ bracket the critical surface allowed by these
residual bounds.

For the conditioning statement, suppose $G\ge2$, $d\ge2$, and
$\widetilde p\ge p_0\in(0,1)$.  Set
\begin{equation}
 q_0=p_0(1-\zeta),\qquad
 \Gamma_*=\lambda_{\max}(C_*-v_*I),\qquad
 \delta_{\max}=\lambda_{\max}(C_L-v_*I).                 \label{eq:cond-robust-constants}
\end{equation}
Then
\begin{equation*}
 x\le\frac{d(1+\eta)\tau}{q_0},                           \condtag{R4}\label{eq:cond-r8}
\end{equation*}
and, whenever the denominator is positive,
\begin{equation*}
 \boxed{\tau\le
 \frac{(1-p_0)(1+\zeta)\delta_{\max}}
 {d\{(1-\eta)(G-1)-(1-p_0)(1+\zeta)(1+\eta)/q_0\}}.}     \condtag{R5}\label{eq:cond-r9}
\end{equation*}
For the log-det Hessian at the approximate point,
\begin{equation}
 \widetilde{\mathcal H}[Y]_g
 =\tau\widetilde Z_g^{-1}Y_g\widetilde Z_g^{-1}
 \quad\text{on }\mathcal T,                               \label{eq:cond-robust-hessian}
\end{equation}
we have
\begin{equation*}
 \boxed{\kappa(\widetilde{\mathcal H}|_{\mathcal T})\ge
 \frac{\delta_{\min}^2(1-\eta)q_0}
 {d(1+\eta)^2\tau(x+\Gamma_*)}=\Omega(G).}                \condtag{R6}\label{eq:cond-r13}
\end{equation*}
If $m\ge2$, the stronger bound is
\begin{equation*}
 \boxed{\kappa(\widetilde{\mathcal H}|_{\mathcal T})\ge
 \frac{\delta_{\min}^2(1-\eta)^2}{(1+\eta)^2x^2}
 =\Omega(G^2).}                                            \condtag{R7}\label{eq:cond-r14}
\end{equation*}
For fixed $p_0,\eta,\zeta,C_*,C_L$, all asymptotic constants are uniform
over such points.
\end{theorem}

\begin{proof}
The residual bound is equivalent to
\[
 (1-\eta)I\preceq\tau^{-1}\widetilde Z_g^{1/2}D_g
 \widetilde Z_g^{1/2}\preceq(1+\eta)I.
\]
Congruence by $\widetilde Z_g^{-1/2}$, inversion, and multiplication by
$\tau$ prove \eqref{eq:cond-r3}.  Since
$D_L\succeq C_L-v_*I$,
\[
 \sum_{g=2}^G\tr\widetilde Z_g
 \le(1+\eta)\tau(G-1)\tr D_L^{-1}
 \le(1+\eta)\tau(G-1)B_0.
\]
Divide by $\widetilde T\ge1-\zeta$ to prove
\eqref{eq:cond-r4}.  The lower sandwich similarly gives
$\widetilde T\ge(1-\eta)\tau(G-1)\tr D_L^{-1}$; taking the limit in the
theorem proves \eqref{eq:cond-r5}.

Winner mass and the upper sandwich give
\[
 q_0\le\tr\widetilde Z_*
 \le(1+\eta)\tau\tr D_*^{-1}
 \le\frac{d(1+\eta)\tau}{x},
\]
which is \eqref{eq:cond-r8}.  Every eigenvalue of $D_L$ is at most
$x+\delta_{\max}$.  The lower sandwich on loser blocks and
$(1-\widetilde p)\widetilde T\le(1-p_0)(1+\zeta)$ give
\[
 (1-p_0)(1+\zeta)\ge
 \frac{(1-\eta)\tau(G-1)d}{x+\delta_{\max}}.
\]
Substitute \eqref{eq:cond-r8} and rearrange to obtain
\eqref{eq:cond-r9}, hence $\tau=O(G^{-1})$.

No commutation is needed for the curvature tests.  In an eigenbasis of any
loser block $\widetilde Z_g$, a Frobenius-unit symmetric off-diagonal matrix
is tangent.  The upper sandwich and $D_L\succeq\delta_{\min}I$ bound every
eigenvalue of this block by $(1+\eta)\tau/\delta_{\min}$, so this direction
has curvature at least
$\delta_{\min}^2/\{(1+\eta)^2\tau\}$.  The largest eigenvalue of the winner
block is at least $q_0/d$, while its other eigenvalues are at least
$(1-\eta)\tau/(x+\Gamma_*)$.  An off-diagonal direction between a largest
eigenvector and an orthogonal eigenvector therefore has curvature at most
$d(x+\Gamma_*)/\{(1-\eta)q_0\}$.  Taking the quotient proves the finite
bound in \eqref{eq:cond-r13}; \eqref{eq:cond-r8}--\eqref{eq:cond-r9} give
its asymptotic form.

If $m\ge2$, the min--max principle and the lower sandwich show that the two
largest winner eigenvalues are at least $(1-\eta)\tau/x$.  Their
off-diagonal direction has curvature at most
$x^2/\{(1-\eta)^2\tau\}$.  Comparing with the loser off-diagonal direction
above proves \eqref{eq:cond-r14}.
\end{proof}

Scalar blocks, excluded above by $d\ge2$, also have a robust linear lower
bound.  When $d=1$ and $G\ge3$, \eqref{eq:cond-r8}--\eqref{eq:cond-r9}
remain valid.
A difference of two loser coordinates has curvature at least
$\delta_1^2/\{(1+\eta)^2\tau\}$.  Balancing a unit winner coordinate by
$-1/(G-1)$ in every loser yields a tangent Rayleigh quotient at most
\begin{equation}
 \frac{\tau}{q_0^2}+
 \frac{(x+\delta_1)^2}{(1-\eta)^2\tau(G-1)}.              \label{eq:cond-robust-scalar-soft}
\end{equation}
The ratio of these two curvature bounds, together with
\eqref{eq:cond-r8}--\eqref{eq:cond-r9}, is $\Omega(G)$.  The residual in
\eqref{eq:cond-robust-residual} must be measured relative to the barrier: a
small unscaled stationarity residual need not imply a Loewner sandwich near
the spectral edge, where $D_*$ becomes singular.

\subsection{Relation to prior work and open questions}
\label{sec:condensation-discussion}

The resolvent form \eqref{eq:cond-exact-center} of the center is classical:
it is the density associated with the matrix Burg entropy, and Ishihara
derives it from a maximum-entropy problem with trace and energy constraints
\cite{ishihara2019-density-operators-generalized-entropies}.  For individual
eigenmodes, the precise statistical-mechanics analogue is Rayleigh--Jeans
condensation, where inverse-gap occupancies, saturation of excited modes,
and a macroscopic ground-mode mass are established
\cite{baudin2020-rayleigh-jeans-condensation,zanaglia2024-bridging-rayleigh-jeans}.
Inverse-gap capacity and condensation onto several states also appear in
spherical models
\cite{lukkarinen2018-multi-state-condensation,crisanti2019-condensation-ordering}.
We therefore call our setting a matrix-Burg / Rayleigh--Jeans analogue and
claim neither a Bose--Einstein distribution nor a new form of the center.

Square-root rates on degenerate SDP central paths and generic ill-conditioning
of Newton systems near low-rank optima are also known
\cite{cruzneto2005-asymptotic-central-path,alizadeh1998-primal-dual-sdp,zhang2017-modified-interior-point-method-for},
but those works do not study the reduced Frobenius Hessian as the number $G$
of blocks grows.  Our contributions are the critical coefficients as $G$
grows, the dependence of the conditioning order ($G$ versus $G^2$) on $m$,
the finite mass--conditioning laws, their combination with visibility and
raw-query bounds, and the conservation law.  The quantum search,
amplification, and adversary tools are standard.

Identical losers are needed for a sharp common capacity, the exact collision
identity, and coherent preparation from public resolvents.  The
average-resolvent hypotheses \eqref{eq:cond-empirical-resolvent} extend the
mass law, but visibility for nonidentical losers requires a separate public
statistic such as \eqref{eq:cond-symmetry-light}.  It is open how the
bound $\eta$ on the barrier-norm residual in \eqref{eq:cond-robust-residual}
relates to the algorithm-specific inexact Newton residuals of
\cref{def:inexact-newton-residual}, whose forcing tolerance is a separate
$\eta$; we claim no conversion.
